\section{Procesamiento con Apache Spark}
\label{sec:processing}

\subsection{Lectura y preparación}

Cada fuente necesitó un tratamiento distinto (Tabla~\ref{tab:preparation}). En todas se renombraron las columnas
al inglés y el resultado se escribió en Parquet.

\begin{table}[htbp]
\centering
\caption{Transformaciones principales de la preparación (notebook 01).}
\label{tab:preparation}
\small
\begin{tabularx}{\textwidth}{@{}l X@{}}
\toprule
Fuente & Transformaciones \\
\midrule
Dengue & Edad normalizada a años (venía en años, meses o días); gravedad a partir del código CIE-10; semana
         epidemiológica convertida en la fecha de su domingo (\code{week\_start}) para alinearla con el clima. \\
RENIEC & Un archivo por año: el año se toma del nombre del archivo con la columna oculta
         \code{\_metadata.file\_path}. Se excluye a los peruanos en el extranjero y se suma la población por
         distrito, sexo y etapa de vida. \\
SUNASS & Archivo en Latin-1; tipo de emergencia (lluvias, déficit hídrico u otra) derivado de la descripción del
         decreto; marca de prórroga. \\
RENIPRESS & Separador \code{;} y direcciones con saltos de línea: sin \code{multiLine=True} y
         \code{escape='"'} se perdían filas. Solo establecimientos activos. \\
SENAMHI & JSON anidado: dos \code{explode} lo llevan a una fila por estación, día y variable, y un \code{pivot}
         devuelve una columna por variable. \\
\bottomrule
\end{tabularx}
\end{table}

\subsection{Análisis}

Las seis preguntas se resuelven con la DataFrame API (notebook 02). Se usan transformaciones (\code{select},
\code{filter}, \code{withColumn}, \code{when}, \code{fillna}, \code{pivot}, \code{unpivot}), agregaciones
(\code{groupBy} y \code{agg} con \code{count}, \code{sum}, \code{avg}, \code{min}, \code{max} y \code{corr}, y
\code{orderBy}) y window functions (\code{lag}, \code{lead}, \code{rowsBetween}, \code{row\_number},
\code{dense\_rank}, \code{ntile}). Cada join se eligió según lo que la pregunta necesita conservar
(Tabla~\ref{tab:joins}).

\begin{table}[htbp]
\centering
\caption{Tipos de join usados y su justificación.}
\label{tab:joins}
\small
\begin{tabularx}{\textwidth}{@{}l l X@{}}
\toprule
Join & Dónde & Justificación \\
\midrule
\code{left} & Q2, Q4, Q5, Q6 & La tabla que guía la consulta conserva todas sus filas: distritos sin
  establecimientos (Q4), distritos con 0 casos (Q5), semanas sin casos (Q6). \\
\code{inner} & Q3 & Casos y población cubren la misma grilla completa (2 sexos × 5 etapas de vida): no se
  pierde nada. \\
\code{left\_anti} & Q2 & Verificación: ningún distrito-año con casos queda sin población (resultado: 0). \\
\code{left\_semi} & Q5 & Filtra los casos de distritos con decreto sin duplicarlos. Un distrito aparece una vez
  por decreto y por prórroga: un inner join devuelve 1,539,044 filas en lugar de 352,213. \\
\code{cross} & Q6 & Cada estación con un calendario de 209 semanas, para que \code{lead} avance semanas y no
  filas cuando faltan datos. \\
\bottomrule
\end{tabularx}
\end{table}

\subsection{DataFrame API y Spark SQL}

La consulta de Q2 (los cinco distritos con mayor incidencia por año) se escribió también en Spark SQL, sobre
temporary views. Las dos versiones devuelven el mismo resultado (\code{exceptAll} en ambas direcciones da 0
filas) y, desde el optimized logical plan, el mismo plan (Sección~\ref{sec:plans}). Por ejemplo, el ranking:

\begin{lstlisting}
# DataFrame API
by_year = Window.partitionBy("year").orderBy(F.desc(F.col("cases") / F.col("population")))
F.dense_rank().over(by_year).alias("rank")
-- Spark SQL
DENSE_RANK() OVER (PARTITION BY year ORDER BY cases / population DESC) AS rank
\end{lstlisting}
